\chapter{Introduction}
\label{ch:introduction}

This chapter introduces our report on Manzil. Manzil is a prototype for smart traffic control and emergency routing that we built at FAST--NUCES. In a city like Lahore, an ambulance often gets stuck in the same queues as every other vehicle. We wanted to see how far signal control, routing and video based counting can help with this problem when they work together. In this chapter we first explain what the report is for. We then explain who we wrote it for and what each group of readers can expect from it. After that we define the terms and abbreviations that we use in the later chapters. One point is important to us throughout the report. We keep a clear line between the work that we have built and measured so far and the work that we still plan to do. We keep that same line in every chapter, so that a reader always knows which results are real and which parts are still only planned.

\section{Purpose of this Document}

This report describes the current state of the Manzil prototype and the evidence that we have collected from it. It covers the requirements, the method and the design of the system, and it also covers how we implemented the system and what we measured. The core of the system is a traffic simulation in SUMO. All of our code talks to SUMO through one simulation adapter, and a safety filter checks every signal change before it is applied. We built three kinds of signal controller. The first is a fixed time controller, and the second is an adaptive controller that uses traffic pressure. The third is a Double DQN controller in which all junctions share the same network. For the ambulance, we built routing that uses travel time and that reroutes the ambulance when a road is closed. We also built a green corridor that moves along ahead of the ambulance. For perception, we trained a YOLO model on Kaggle and used it in a pipeline that detects, tracks and counts vehicles. Lastly, we built a first version of a web dashboard.

All of our controller, routing and corridor results come from repeated SUMO runs with fixed seeds. These runs use a synthetic 3$\times$3 grid of signalized junctions, and they do not use a real city map yet. All of our perception results come from the public UA-DETRAC and COCO datasets. So the evidence in this report is measured evidence from simulation and from video. We do not claim how the system would perform on real roads. We also do not claim that it is ready to deploy. Whenever we make a statement about performance in this report, we tie it to a recorded configuration and a matched baseline. We also tie it to a stated metric and to a set of runs that anyone can repeat. In this way a reader can check each number in the report against the files in our project repository, and the reader can also run the same experiment again to get the same result.

\section{Intended Audience}

We wrote this report mainly for the academic readers at FAST--NUCES who will judge the project. These readers need enough detail to judge whether our objectives make sense, what software we built and whether our experiments are valid. The report is also an engineering reference for us, the three members of the team. From the chapters that follow, each of us should be able to see what every component is responsible for and which fields each interface carries. Each of us should also be able to find how each scenario is set up, how the system falls back when something fails and how each part is tested. There is also a second group of readers. These are people in traffic management and emergency services who may want to understand the problem, our assumptions and the limits of the prototype. We do not present them as our clients, our partners or our approving authorities. For this reason, we explain each domain term when it first becomes important, but we still keep the detail that is needed for implementation and academic review. We keep business points apart from technical evidence, and we do not describe a proposed function as if it already works. Nobody should mistake our research prototype for a traffic control system that is in real use.

\section{Definitions, Acronyms, and Abbreviations}

Our project mixes ideas from traffic engineering, computer vision and reinforcement learning, so the same word can mean different things to different readers. For example, a queue in SUMO is an exact number that the simulator gives us. A queue from video is only an estimate that we get from tracked vehicles. In the same way, a green corridor in our work is a set of limited and timed signal changes that move along with the ambulance, and the signals on the route still follow their safety rules. To avoid this kind of confusion, we define below the main terms and abbreviations that we use in this report. We give a short meaning for each term. Where a term comes from a specific paper or tool, we also cite that paper or tool. We use these meanings in the same way in every later chapter of the report, so the reader can come back to this list whenever a term is not clear.

\begin{description}[font=\normalfont,leftmargin=0pt,itemindent=0.45em,labelsep=0.45em,itemsep=0pt,parsep=0pt,topsep=0.25\baselineskip]
\item[A*:] A shortest path search. It adds up the cost so far and uses a lower bound heuristic that never overestimates the remaining cost \cite{hart1968astar}.
\item[API:] Application Programming Interface.
\item[ByteTrack:] A tracker for multiple objects. It first matches the detections with high confidence. It then uses the detections with lower confidence to recover objects that it would otherwise lose \cite{zhang2022bytetrack}.
\item[CV:] Computer vision. In this report we use it to get traffic indicators that can be observed from video.
\item[DQN:] Deep Q Network. It is a reinforcement learning method that uses a neural network to estimate the value of each action.
\item[Double DQN:] A version of DQN that picks the next action with the online network and then scores that action with a target network. This reduces overestimation \cite{vanhasselt2016double}.
\item[Green corridor:] A sequence of limited signal priority changes that moves along ahead of an ambulance.
\item[MAE:] Mean absolute error. It is the average absolute difference between the estimated values and the true values.
\item[mAP50:] Mean average precision when the intersection over union threshold is 0.5. It is a standard measure of detection accuracy.
\item[MOT:] Multiple object tracking. It links detections across video frames so that each object keeps one identity.
\item[Queue proxy:] An estimate of the number of tracked vehicles that are slow or stopped in a defined lane region.
\item[RL:] Reinforcement learning. In RL a policy is improved using states, actions and the total reward over time.
\item[Signal preemption:] A controlled change of signal for an emergency request. It includes the change itself and the recovery after it.
\item[SUMO:] Simulation of Urban MObility. It is a microscopic road traffic simulator.
\item[Throughput:] The number of vehicles that leave a measured movement or network region in a unit of time.
\item[TraCI:] Traffic Control Interface. It lets a program read and change a SUMO simulation while the simulation runs. libsumo offers the same interface inside the Python process.
\item[Traffic state contract:] A versioned structure for observations. It records the source, time, units, confidence and validity of each observation.
\item[UA-DETRAC:] A public benchmark for detecting and tracking vehicles in video from fixed traffic cameras \cite{wen2020uadetrac}.
\item[Waiting time:] The time that a vehicle spends below the speed threshold that we set in SUMO.
\item[WAPE:] Weighted absolute percentage error. It is the sum of the absolute count errors divided by the sum of the true counts.
\item[YOLO:] You Only Look Once. It is a family of single stage object detection methods.
\item[YOLO11s:] The small model in the Ultralytics YOLO11 detector family \cite{ultralytics2024yolo11}.
\end{description}

\section{Conclusion}

In this chapter we presented Manzil as one study of several traffic decisions that affect each other, and we built it so that others can repeat our work. The system looks at how observed traffic conditions can guide signal control and how routing can react to delay in the network. It also looks at how an ambulance can get priority for a short time while normal traffic still keeps moving. We also set the limits of our evidence in this chapter. Our repository now contains working controllers, a trained detector, a trained DQN policy, a synthetic SUMO network that we have validated and the recorded results of our experiments. We report all of these as measured evidence. Some work is still left for the next phase of the project. This includes the Lahore road network, the full dashboard with five screens and a database, video that we film locally and the full protocol with ten seeds. We describe all of this as planned work. In this report we call something implemented only if its code, configuration, test evidence or experiment record exists in our project repository.

The rest of the report is organized as follows. Chapter \ref{ch:vision} presents the vision, objectives and scope of the project. Chapter \ref{ch:literature} reviews related research on traffic perception, signal control, reinforcement learning and emergency routing, and it shows where our work fits. Chapter \ref{ch:srs} gives the functional and nonfunctional requirements of the system. Chapter \ref{ch:method} explains our approach and our experimental protocol, including the baselines and the metrics that we use. Chapter \ref{ch:design} gives the high level design and the low level design of the system. Chapter \ref{ch:implementation} describes the implementation and the test cases that we wrote to check it. Chapter \ref{ch:manual} is a short user manual for the prototype, and it shows how to run the simulation and the dashboard. Chapter \ref{ch:results} reports our experimental results and discusses what they mean. Finally, Chapter \ref{ch:conclusion} concludes the report and gives our plan for the next phase of the work.
